\subchapter{A/B updates}{Objective:
	Learn to configure RAUC and deploy A/B updates to our board}

\section{Prerequisites}

For this stage, we are going to run RAUC also on the development host, to generate
the update bundle(s).
Run the following command:

\begin{bashinput}
$ sudo apt install -y rauc squashfs-tools
\end{bashinput}

\section{A/B update partition setup}

To avoid rebuilds and reflashes, the system already has an A/B setup.
You can confirm this and look at the details in \code{board_setup/meta-security-training/wic/security-training-image.freiheit93.wks.in}

\section{Configuring RAUC on the platform}

Systemd should have tried to start the RAUC daemon:

\begin{bashinput}
# systemctl status rauc
\end{bashinput}

This should show that the service has not been started.
You can look at the logs to figure out why:

\begin{bashinput}
# journalctl -xeu rauc.service
\end{bashinput}

RAUC needs a few things it currently does not have to function:
\begin{itemize}
  \item its main \href{https://rauc.readthedocs.io/en/latest/reference.html\#system-configuration-file}{configuration file}
  \item the \href{https://rauc.readthedocs.io/en/latest/reference.html\#keyring-section}{keyring} which will be used
	  to verify the bundle's signature
  \item functional \code{fw_printenv} and \code{fw_setenv} commands
\end{itemize}

\subsection{Setting up fw\_printenv/fw\_setenv}

Running \code{fw_printenv} should currently fail:

\begin{bashinput}
$ fw_printenv
Cannot initialize environment
\end{bashinput}

The reason for this is that the \href{https://github.com/ARM-software/u-boot/blob/master/tools/env/fw_env.config}{configuration} is missing.
Considering the A/B setup as defined in \code{board_setup/meta-security-training/wic/security-training-image.freiheit93.wks.in}, create the
configuration file in \code{/etc/fw_env.config}.


\subsection{The \href{https://rauc.readthedocs.io/en/latest/reference.html\#keyring-section}{keyring}}

RAUC does not support \textbf{unsigned} updates, by design. Therefore, we need a key pair for our RAUC updates.
We'll also want a certificate, that we'll deploy to the target for RAUC to use to verify the signature.
We could use:
\begin{itemize}
  \item A self-signed certificate
  \item the root CA certificate
\end{itemize}

The first will mean less flexibility in handling certificates, while the second risks exposing
your root key, which is unnecessary.
Fortunately, we already have the key pair (and certificate) we need: the bundle will be built on the host,
we'll just use its key pair and certificate.

We could use \code{host.crt} directly as the keyring, but using \code{cacert.pem} is better, as we expect
that certificate to change very infrequently, and it is the root of trust of our PKI.

In a real case, we would possibly have at least one child CA, and therefore one intermediate certificate,
but we'll keep things simple here.

\subsection{Building and installing the bundle}

RAUC bundles use a configuration file called a \href{https://rauc.readthedocs.io/en/latest/reference.html\#sec-ref-manifest}{manifest}.
This is what you'll be writing in this section.

Once you're done, you can build the update bundle:

\begin{bashinput}
$ mkdir bundle
# edit bundle/manifest.raucm
$ cd bundle
$ mkdir root
$ echo "The rootfs has been updated" > root/UPDATED.txt
$ gzip -d rootfs.tar.gz
$ tar -u rootfs.tar root
$ rauc bundle . update.bundle --signing-keyring=../ca/cacert.pem --key=../host/hostkey.pem --cert=../host/host.crt 
Creating 'verity' format bundle
rauc-Message: Keyring given, doing signature verification
rauc-Message: Verified inline signature by 'C = FR, ST = Some-State, O = Bootlin security training, CN = host computer'
$ scp update.bundle root@{BOARD}
\end{bashinput}

And then install it and reboot:

\begin{bashinput}
$ rauc install update.bundle
\end{bashinput}

RAUC might fail because the system's date is too long in the past for the certificate we generated to already be valid.

\begin{bashinput}
$ date -s 'YYYY-MM-DD'
\end{bashinput}

Alternatively, you can also instruct RAUC to use the time at which the bundle was signed instead of the system time,
by adding
\begin{bashinput}
use-bundle-signing-time=true
\end{bashinput}
to the \code{[keyring]} section of \code{/etc/rauc/system.conf}.

Even once the bundle is correctly installed, the board should not show UPDATED.txt.
Running \code{mount} should help figure out why.

You can also look at the U-Boot environment to see what RAUC attempted to do:
\begin{bashinput}
$ fw_printenv
BOOT_ORDER=rootB
BOOT_rootB_LEFT=3
\end{bashinput}

\subsection{U-Boot configuration}

We now need to change the way U-Boot boots up the system.
You can take some inspiration from \href{https://github.com/rauc/rauc/blob/master/contrib/uboot.sh}{this example}.
